Penelitian ini mengimplementasikan dan menguji kontrol KYC dan AML/CFT terpilih
pada purwarupa Rupiah Digital ritel berbasis Hyperledger Fabric 2.5. Lingkungan
studi kontrol menggunakan jaringan terisolasi dengan empat MSP:
\texttt{CentralBankMSP}, \texttt{DistributorMSP}, \texttt{ValidatorMSP}, dan
\texttt{EndUserMSP}. Tiga \textit{chaincode} Go memisahkan operasi moneter,
kontrol pelanggan, dan pembayaran. Identitas X.509 yang dikelola \textit{backend}
memisahkan penanda tangan, subjek pelanggan, dan kustodian dompet.

Uji unit sintetis dan uji integrasi \textit{chaincode} lulus untuk keputusan KYC
berisiko tinggi dengan peninjau dan pejabat senior yang berbeda, pembatasan dan
pembukaan pembatasan berbukti, pendanaan cadangan, pembayaran ritel, pembacaan
oleh \texttt{ValidatorMSP}, serta penolakan transaksi dengan \textit{endorsement}
yang tidak mencukupi. Hasil ini hanya menunjukkan perilaku pada skenario sintetis
yang diuji; penelitian tidak membuktikan keaslian identitas, efektivitas
penyaringan, pengiriman laporan kepada PPATK, atau kepatuhan institusional.

Sebagai bukti kinerja sekunder, Hyperledger Caliper menguji baseline terpisah
berisi lima organisasi dengan \texttt{digital-rupiah} v3.0. Transfer ritel pada
kontensi rendah mencapai 20 TPS, sedangkan uji berpasangan pada kontensi tinggi
mencapai 2,01 TPS. Pengukuran tersebut tidak menguji tiga \textit{chaincode}
lingkungan kontrol dan tidak menetapkan kapasitas produksi atau reproduktibilitas
lintas mesin.

\vspace{0.5cm}

\noindent{Kata kunci}: Rupiah Digital, CBDC ritel, KYC, AML/CFT,
Hyperledger Fabric, \textit{endorsement}, distribusi dua jenjang
